\documentclass[11pt,a4paper]{article}

% ============================================================
% LIVRET INTENSIF — ÉPREUVE ANTICIPÉE DE MATHÉMATIQUES
% PREMIÈRE SPÉCIALITÉ MATHÉMATIQUES
% Version longue structurée, prête à compiler.
% ============================================================

\usepackage[utf8]{inputenc}
\usepackage[T1]{fontenc}
\usepackage[french]{babel}
\usepackage{lmodern}
\usepackage{geometry}
\geometry{margin=1.65cm}

\usepackage{amsmath,amssymb,amsthm}
\usepackage{array}
\usepackage{tabularx}
\usepackage{longtable}
\usepackage{booktabs}
\usepackage{multicol}
\usepackage{xcolor}
\usepackage{enumitem}
\usepackage{fancyhdr}
\usepackage{hyperref}
\usepackage[most]{tcolorbox}
\usepackage{pifont}

\hypersetup{
  colorlinks=true,
  linkcolor=blue!60!black,
  urlcolor=blue!60!black
}

\pagestyle{fancy}
\fancyhf{}
\lhead{Programme intensif — Première spé maths}
\rhead{Épreuve anticipée de mathématiques}
\cfoot{\thepage}

\setlist[itemize]{leftmargin=1.25em}
\setlist[enumerate]{leftmargin=1.55em}

\definecolor{bleuprepa}{RGB}{20,55,110}
\definecolor{vertprepa}{RGB}{20,105,70}
\definecolor{rougeprepa}{RGB}{160,35,35}
\definecolor{orangeprepa}{RGB}{180,95,20}
\definecolor{violetprepa}{RGB}{95,45,130}
\definecolor{grisclair}{RGB}{245,247,250}

\newcommand{\casevide}{\ding{113}}
\newcommand{\important}[1]{\textbf{\textcolor{rougeprepa}{#1}}}
\newcommand{\sujetlink}[2]{\hyperlink{#1}{\textcolor{blue!60!black}{\textbf{#2}}}}

\newcommand{\R}{\mathbb{R}}
\newcommand{\N}{\mathbb{N}}
\newcommand{\e}{\mathrm{e}}

\tcbset{
  enhanced,
  breakable,
  sharp corners=downhill,
  boxrule=0.7pt,
  colback=white,
  fonttitle=\bfseries
}

\newtcolorbox{objectifjour}[1][]{
  colframe=bleuprepa,
  colback=blue!2!white,
  title={Objectif du jour #1}
}

\newtcolorbox{methode}[1][]{
  colframe=vertprepa,
  colback=green!2!white,
  title={Méthode #1}
}

\newtcolorbox{alerte}[1][]{
  colframe=rougeprepa,
  colback=red!2!white,
  title={Alerte #1}
}

\newtcolorbox{sujetofficiel}[1][]{
  colframe=bleuprepa,
  colback=cyan!2!white,
  title={Sujet officiel ou ressource à faire #1}
}

\newtcolorbox{sujetprof}[1][]{
  colframe=violetprepa,
  colback=violetprepa!3!white,
  title={Sujet de professeur / ressource pédagogique #1}
}

\newtcolorbox{sujettype}[1][]{
  colframe=bleuprepa,
  colback=blue!1!white,
  title={Sujet type épreuve inventé #1}
}

\newtcolorbox{difficile}[1][]{
  colframe=orangeprepa,
  colback=orange!3!white,
  title={Sujet plus difficile que l'épreuve #1}
}

\newtcolorbox{corrige}[1][]{
  colframe=vertprepa,
  colback=green!1!white,
  title={Correction synthétique #1}
}

\title{
  \Huge\textbf{Programme intensif de révision}\\[2mm]
  \Large Première spécialité mathématiques\\[1mm]
  \large Épreuve anticipée de mathématiques — objectif excellence
}
\author{Livret professeur — Première générale}
\date{Version prête à compiler}

\begin{document}

\maketitle

\begin{alerte}[Important — esprit du livret]
Ce livret ne recopie pas mot pour mot les sujets existants. Il donne des références précises et des liens cliquables vers des ressources officielles, des sujets zéro, des annales, des devoirs et des exercices de professeurs. Il propose ensuite des sujets originaux inventés dans l'esprit de l'épreuve.
\end{alerte}

\begin{alerte}[Format de l'épreuve]
L'épreuve anticipée de mathématiques pour les élèves de première spécialité mathématiques dure \textbf{2 heures}.  
Elle est conçue sans calculatrice et comporte :
\begin{itemize}
  \item une partie d'automatismes, souvent sous forme de QCM ;
  \item une partie d'exercices rédigés portant sur le programme de première spécialité mathématiques.
\end{itemize}
\end{alerte}

\tableofcontents
\newpage

% ============================================================
\section{Sources officielles et liens utiles}
% ============================================================

\begin{sujetofficiel}[sujets zéro officiels]
\begin{itemize}
  \item \href{https://eduscol.education.fr/sites/default/files/document/sujet-specialite-1pdf-112050.pdf}{Eduscol — Sujet zéro officiel — Voie générale, spécialité mathématiques}.
  \item \href{https://www.apmep.fr/Epreuve-anticipee-2025}{APMEP — Épreuve anticipée 2026 et sujets zéro}.
  \item \href{https://math93.com/index.php?Itemid=1608&catid=164%3Aannales-du-bac-premiere-maths-epreuve-anticipee-de-mathematiquesde-premiere&id=1192%3Abac-2026-epreuve-anticipee-de-mathematiques-en-premiere-sujets-et-corriges-math93&option=com_content&view=article}{Math93 — Épreuve anticipée de mathématiques en première, sujets et corrigés}.
\end{itemize}
\end{sujetofficiel}

\begin{sujetprof}[banques de devoirs de professeurs]
\begin{itemize}
  \item \href{https://www.lyceedadultes.fr/sitepedagogique/pages/math1Spe_controlesdevoirs.html}{Lycée d'adultes — Contrôles et devoirs de première spécialité}.
  \item \href{https://www.annales2maths.com/exercices-corriges-1s/}{Annales2Maths — Exercices corrigés de première}.
  \item \href{https://mangeard.maths.free.fr/premiereS.htm}{Mangeard Maths — Contrôles, corrigés et exercices de première}.
  \item \href{https://mathexien.com/Maths1ere/assets/pdf/ds_commun_2022_1_correction.pdf}{MatheX — Correction devoir commun de première spé maths}.
  \item \href{https://perpendiculaires.free.fr/?page_id=2554}{Perpendiculaires — Polycopiés de première générale spécialité mathématiques}.
  \item \href{https://lionelponton.fr/Premi%C3%A8re-sp%C3%A9cialit%C3%A9/}{Lionel Ponton — Première spécialité, cours, exercices, interrogations}.
  \item \href{https://www.maths-et-tiques.fr/index.php/cours-maths/niveau-premiere}{Maths et tiques — Cours et exercices niveau première}.
  \item \href{https://cours-galilee.com/ressources-2/specialite-maths-premiere/exercices-maths-premiere/}{Cours Galilée — Exercices corrigés de première spécialité}.
  \item \href{https://maths-pdf.fr/exercices-maths-1ere}{Maths-PDF — Exercices corrigés de première}.
\end{itemize}
\end{sujetprof}

\begin{alerte}[Utilisation des ressources externes]
Les liens vers des devoirs de professeurs doivent être utilisés comme \textbf{banque d'entraînement}.  
On ne les recopie pas dans ce livret : on les fait en temps limité, puis on corrige activement les erreurs.
\end{alerte}

% ============================================================
\section{Objectif général de la préparation}
% ============================================================

\subsection{Ce que signifie réussir l'épreuve}

Réussir l'épreuve anticipée de mathématiques ne signifie pas seulement connaître les formules.  
Il faut être capable de :
\begin{itemize}
  \item calculer vite et juste sans calculatrice ;
  \item reconnaître immédiatement le chapitre mobilisé ;
  \item rédiger proprement les exercices ;
  \item éviter les erreurs de signe ;
  \item justifier les étapes importantes ;
  \item gérer les 2 heures ;
  \item ne pas bloquer sur une question ;
  \item exploiter efficacement les questions précédentes.
\end{itemize}

\begin{methode}[stratégie en 2 heures]
\begin{enumerate}
  \item Lire tout le sujet en 3 à 5 minutes.
  \item Commencer par les automatismes.
  \item Ne pas passer trop de temps sur une question de QCM.
  \item Pour les exercices, identifier le chapitre : second degré, dérivation, suites, probabilités, géométrie.
  \item Rédiger les questions faciles proprement.
  \item Laisser une trace même si une question est difficile.
  \item Garder 10 à 15 minutes pour relire les calculs.
\end{enumerate}
\end{methode}

\begin{methode}[relecture mathématique]
\begin{multicols}{2}
\begin{itemize}
  \item Fractions simplifiées.
  \item Signe des expressions.
  \item Domaine de définition.
  \item Racines du trinôme.
  \item Tableau de signes.
  \item Tableau de variations.
  \item Dérivée correcte.
  \item Tangente bien écrite.
  \item Probabilité entre 0 et 1.
  \item Somme des probabilités égale à 1.
  \item Coordonnées vérifiées.
  \item Conclusion avec phrase.
\end{itemize}
\end{multicols}
\end{methode}

% ============================================================
\section{Hiérarchie des chapitres}
% ============================================================

\footnotesize
\begin{longtable}{|p{3.2cm}|p{2.4cm}|p{3.2cm}|p{4cm}|p{4cm}|}
\hline
\textbf{Chapitre} & \textbf{Priorité} & \textbf{Type de questions} & \textbf{Méthodes indispensables} & \textbf{Niveau excellence}\\
\hline
Calcul algébrique & Absolue & Automatismes, QCM & Développer, factoriser, simplifier & Zéro erreur de signe ou de fraction.\\
\hline
Second degré & Absolue & Équation, signe, parabole & Discriminant, racines, forme canonique & Savoir choisir la méthode la plus rapide.\\
\hline
Fonctions & Absolue & Lecture graphique, variations & Images, antécédents, tableau de variations & Interpréter graphiquement et algébriquement.\\
\hline
Dérivation & Absolue & Tangente, variations, extremum & Dérivée, signe de la dérivée, tableau & Rédaction propre et conclusion claire.\\
\hline
Suites & Très important & Modélisation, explicite, récurrence & Calcul de termes, variations, seuil & Savoir passer d'une situation concrète à une suite.\\
\hline
Probabilités & Très important & Arbres, conditionnement & Probabilités totales, contraire, intersection & Distinguer $P(A\cap B)$ et $P_A(B)$.\\
\hline
Géométrie repérée & Important & Vecteurs, droites, produit scalaire & Coordonnées, colinéarité, orthogonalité & Rédiger une démonstration géométrique complète.\\
\hline
Trigonométrie & Important & Cercle trigonométrique, équations simples & Valeurs remarquables, sinus, cosinus & Aller vite sur les valeurs classiques.\\
\hline
Algorithmique & Important & Suites, seuils, boucles & Lire et écrire un algorithme simple & Expliquer le rôle d'une boucle.\\
\hline
\end{longtable}
\normalsize

% ============================================================
\section{Planning intensif avec liens cliquables vers sujets de profs}
% ============================================================

\begin{alerte}[principe du planning]
Le planning est écrit en jours relatifs : J-21, J-20, etc.  
Il peut donc être utilisé quelle que soit la date de départ.  
Chaque séance contient :
\begin{itemize}
  \item une révision ciblée ;
  \item un lien vers un sujet ou une ressource existante ;
  \item un sujet inventé interne au livret ;
  \item une case à cocher.
\end{itemize}
\end{alerte}

\footnotesize
\begin{longtable}{|p{1.5cm}|p{2.6cm}|p{6.6cm}|p{5.6cm}|p{1.1cm}|}
\hline
\textbf{Jour} & \textbf{Objectif} & \textbf{Travail principal} & \textbf{Ressource externe à utiliser} & \textbf{Fait}\\
\hline

J-21 &
Diagnostic complet &
Faire le sujet inventé \sujetlink{P-SUJET-01}{P-SUJET-01 — Diagnostic première spé}.  
Durée : 2 h.  
Noter les chapitres faibles. &
Lire le sujet zéro officiel :  
\href{https://eduscol.education.fr/sites/default/files/document/sujet-specialite-1pdf-112050.pdf}{Eduscol — Sujet zéro spécialité mathématiques}. &
\casevide\\
\hline

J-20 &
Automatismes sans calculatrice &
Faire 45 minutes de calculs : fractions, puissances, racines, identités remarquables, équations simples.  
Puis faire \sujetlink{P-AUTO-01}{P-AUTO-01 — Automatismes niveau épreuve}. &
Utiliser :
\href{https://www.maths-et-tiques.fr/index.php/cours-maths/niveau-premiere}{Maths et tiques — Niveau première}.  
Choisir les exercices courts de calcul. &
\casevide\\
\hline

J-19 &
Second degré I &
Réviser discriminant, racines, forme canonique.  
Faire \sujetlink{P-SUJET-02}{P-SUJET-02 — Second degré et paraboles}. &
Faire un devoir ou exercice du chapitre second degré sur :
\href{https://www.lyceedadultes.fr/sitepedagogique/pages/math1Spe_controlesdevoirs.html}{Lycée d'adultes — Contrôles et devoirs 1ère spé}. &
\casevide\\
\hline

J-18 &
Second degré II &
Étudier le signe d'un trinôme et résoudre des inéquations.  
Refaire les erreurs de J-19. &
Utiliser les exercices corrigés de :
\href{https://www.annales2maths.com/exercices-corriges-1s/}{Annales2Maths — Second degré, dérivation, suites}. &
\casevide\\
\hline

J-17 &
Fonctions générales &
Images, antécédents, lectures graphiques, variations.  
Faire \sujetlink{P-SUJET-03}{P-SUJET-03 — Fonction et lecture graphique}. &
Utiliser :
\href{https://mangeard.maths.free.fr/premiereS.htm}{Mangeard Maths — Contrôles et corrigés de première}. &
\casevide\\
\hline

J-16 &
Dérivation I &
Nombre dérivé, tangente, équation de tangente.  
Faire \sujetlink{P-SUJET-04}{P-SUJET-04 — Dérivation et tangentes}. &
Consulter :
\href{https://www.maths-et-tiques.fr/index.php/cours-maths/niveau-premiere}{Maths et tiques — Dérivation en première}. &
\casevide\\
\hline

J-15 &
Dérivation II &
Calcul de dérivées et signe de la dérivée.  
Faire une étude complète de fonction. &
Faire un devoir sur la dérivation depuis :
\href{https://www.lyceedadultes.fr/sitepedagogique/pages/math1Spe_controlesdevoirs.html}{Lycée d'adultes — Devoirs première spé}. &
\casevide\\
\hline

J-14 &
Sujet complet 1 &
Faire \sujetlink{P-SUJET-05}{P-SUJET-05 — Sujet complet type épreuve 1}.  
Durée : 2 h.  
Correction active : 1 h. &
Comparer avec :
\href{https://www.apmep.fr/Epreuve-anticipee-2025}{APMEP — Sujets zéro et corrigés}. &
\casevide\\
\hline

J-13 &
Suites I &
Suites explicites, suites récurrentes, premiers termes.  
Faire \sujetlink{P-SUJET-06}{P-SUJET-06 — Suites et modélisation}. &
Utiliser :
\href{https://perpendiculaires.free.fr/?page_id=2554}{Perpendiculaires — Polycopiés première spécialité}. &
\casevide\\
\hline

J-12 &
Suites II &
Suites arithmétiques et géométriques.  
Seuils et algorithmes. &
Utiliser :
\href{https://lionelponton.fr/Premi%C3%A8re-sp%C3%A9cialit%C3%A9/}{Lionel Ponton — Cours et exercices première spécialité}. &
\casevide\\
\hline

J-11 &
Probabilités I &
Vocabulaire, événements, contraire, union, intersection.  
Faire \sujetlink{P-SUJET-07}{P-SUJET-07 — Probabilités dans un lycée}. &
Faire une fiche probabilités depuis :
\href{https://www.annales2maths.com/exercices-corriges-1s/}{Annales2Maths — Probabilités première}. &
\casevide\\
\hline

J-10 &
Probabilités II &
Arbres pondérés, probabilités conditionnelles.  
Faire \sujetlink{P-SUJET-08}{P-SUJET-08 — Test médical et probabilités conditionnelles}. &
Utiliser :
\href{https://lionelponton.fr/Premi%C3%A8re-sp%C3%A9cialit%C3%A9/}{Lionel Ponton — Probabilités conditionnelles}. &
\casevide\\
\hline

J-9 &
Géométrie repérée I &
Coordonnées, vecteurs, milieu, distance.  
Faire \sujetlink{P-SUJET-09}{P-SUJET-09 — Géométrie repérée}. &
Utiliser les polycopiés de :
\href{https://perpendiculaires.free.fr/?page_id=2554}{Perpendiculaires — Géométrie première}. &
\casevide\\
\hline

J-8 &
Géométrie repérée II &
Droites, colinéarité, produit scalaire, orthogonalité. &
Chercher un contrôle de géométrie dans :
\href{https://mangeard.maths.free.fr/premiereS.htm}{Mangeard Maths — Première}. &
\casevide\\
\hline

J-7 &
Sujet complet 2 &
Faire \sujetlink{P-SUJET-10}{P-SUJET-10 — Sujet complet type épreuve 2}.  
Durée : 2 h sans calculatrice. &
Faire un sujet zéro depuis :
\href{https://math93.com/index.php?Itemid=1608&catid=164%3Aannales-du-bac-premiere-maths-epreuve-anticipee-de-mathematiquesde-premiere&id=1192%3Abac-2026-epreuve-anticipee-de-mathematiques-en-premiere-sujets-et-corriges-math93&option=com_content&view=article}{Math93 — Sujets zéro et corrigés}. &
\casevide\\
\hline

J-6 &
Trigonométrie &
Cercle trigonométrique, valeurs remarquables, équations simples.  
Faire \sujetlink{P-SUJET-11}{P-SUJET-11 — Trigonométrie}. &
Utiliser :
\href{https://www.maths-et-tiques.fr/index.php/cours-maths/niveau-premiere}{Maths et tiques — Trigonométrie première}. &
\casevide\\
\hline

J-5 &
Algorithmique &
Boucles, seuils, suites, fonctions Python simples.  
Faire \sujetlink{P-SUJET-12}{P-SUJET-12 — Algorithmique et suites}. &
Utiliser des exercices de suites et algorithmes sur :
\href{https://www.annales2maths.com/exercices-corriges-1s/}{Annales2Maths — Exercices corrigés première}. &
\casevide\\
\hline

J-4 &
Sujet complet 3 &
Faire \sujetlink{P-SUJET-13}{P-SUJET-13 — Sujet complet type épreuve 3}.  
Correction immédiate. &
Faire un devoir commun depuis :
\href{https://mathexien.com/Maths1ere/assets/pdf/ds_commun_2022_1_correction.pdf}{MatheX — Devoir commun corrigé}. &
\casevide\\
\hline

J-3 &
Consolidation ciblée &
Refaire uniquement les exercices ratés.  
Remplir le carnet d'erreurs. &
Reprendre les liens déjà utilisés :
\href{https://www.lyceedadultes.fr/sitepedagogique/pages/math1Spe_controlesdevoirs.html}{Lycée d'adultes},
\href{https://www.annales2maths.com/exercices-corriges-1s/}{Annales2Maths},
\href{https://perpendiculaires.free.fr/?page_id=2554}{Perpendiculaires}. &
\casevide\\
\hline

J-2 &
Dernier sujet léger &
Faire seulement les automatismes et un exercice de fonction.  
Pas de gros nouveau chapitre. &
Relire :
\href{https://eduscol.education.fr/sites/default/files/document/sujet-specialite-1pdf-112050.pdf}{Sujet zéro officiel Eduscol}. &
\casevide\\
\hline

J-1 &
Repos intelligent &
Relire les fiches méthodes.  
Refaire 10 automatismes.  
Préparer le matériel. &
Aucune nouvelle ressource.  
On relit uniquement ses erreurs personnelles. &
\casevide\\
\hline

Jour J &
Épreuve &
Lire tout le sujet.  
Commencer par les questions rentables.  
Garder 10 à 15 minutes de relecture. &
Ne pas paniquer : une question difficile peut souvent être sautée puis reprise. &
\casevide\\
\hline

\end{longtable}
\normalsize

% ============================================================
\section{Fiches méthodes indispensables}
% ============================================================

\begin{methode}[calculer sans calculatrice]
\begin{enumerate}
  \item Simplifier les fractions avant de multiplier.
  \item Transformer une soustraction en addition de l'opposé si nécessaire.
  \item Vérifier les signes à chaque ligne.
  \item Utiliser les identités remarquables dès qu'elles apparaissent.
  \item Ne jamais développer inutilement si une factorisation est plus efficace.
\end{enumerate}
\end{methode}

\begin{methode}[résoudre une équation du second degré]
Pour résoudre
\[
ax^2+bx+c=0,\qquad a\neq0,
\]
on calcule
\[
\Delta=b^2-4ac.
\]
\begin{itemize}
  \item Si $\Delta>0$, deux racines :
  \[
  x_1=\frac{-b-\sqrt{\Delta}}{2a},
  \qquad
  x_2=\frac{-b+\sqrt{\Delta}}{2a}.
  \]
  \item Si $\Delta=0$, une racine double :
  \[
  x_0=\frac{-b}{2a}.
  \]
  \item Si $\Delta<0$, aucune racine réelle.
\end{itemize}
\end{methode}

\begin{methode}[étudier le signe d'un trinôme]
Pour étudier le signe de $ax^2+bx+c$ :
\begin{enumerate}
  \item calculer le discriminant ;
  \item trouver les racines si elles existent ;
  \item placer les racines dans un tableau de signes ;
  \item utiliser le signe de $a$ ;
  \item conclure avec des intervalles.
\end{enumerate}
Quand $\Delta>0$, le trinôme est du signe de $a$ à l'extérieur des racines.
\end{methode}

\begin{methode}[étudier une fonction avec la dérivée]
\begin{enumerate}
  \item Déterminer l'ensemble de définition.
  \item Calculer la dérivée.
  \item Factoriser la dérivée si possible.
  \item Étudier le signe de la dérivée.
  \item Dresser le tableau de variations.
  \item Répondre précisément à la question posée.
\end{enumerate}
\end{methode}

\begin{methode}[équation de tangente]
L'équation de la tangente à la courbe de $f$ au point d'abscisse $a$ est :
\[
y=f'(a)(x-a)+f(a).
\]
Il faut donc toujours calculer :
\[
f(a)\quad\text{et}\quad f'(a).
\]
\end{methode}

\begin{methode}[suites]
Pour une suite, commencer par identifier le type de définition.
\begin{itemize}
  \item Suite explicite :
  \[
  u_n=f(n).
  \]
  \item Suite récurrente :
  \[
  u_{n+1}=f(u_n).
  \]
  \item Suite arithmétique :
  \[
  u_{n+1}=u_n+r,
  \qquad
  u_n=u_0+nr.
  \]
  \item Suite géométrique :
  \[
  u_{n+1}=qu_n,
  \qquad
  u_n=u_0q^n.
  \]
\end{itemize}
\end{methode}

\begin{methode}[probabilités conditionnelles]
Si $P(A)\neq0$, alors :
\[
P_A(B)=\frac{P(A\cap B)}{P(A)}.
\]
Donc :
\[
P(A\cap B)=P(A)\times P_A(B).
\]
Dans un arbre pondéré, on multiplie le long d'une branche et on additionne plusieurs chemins incompatibles.
\end{methode}

\begin{methode}[géométrie repérée]
Si $A(x_A;y_A)$ et $B(x_B;y_B)$, alors :
\[
\overrightarrow{AB}(x_B-x_A\,;\,y_B-y_A).
\]
La distance est :
\[
AB=\sqrt{(x_B-x_A)^2+(y_B-y_A)^2}.
\]
Le milieu de $[AB]$ est :
\[
I\left(\frac{x_A+x_B}{2};\frac{y_A+y_B}{2}\right).
\]
\end{methode}

\begin{methode}[produit scalaire]
Si $\vec u(x;y)$ et $\vec v(x';y')$, alors :
\[
\vec u\cdot \vec v=xx'+yy'.
\]
Deux vecteurs sont orthogonaux si et seulement si :
\[
\vec u\cdot \vec v=0.
\]
\end{methode}

% ============================================================
\section{Sujets inventés pour l'épreuve}
% ============================================================

\hypertarget{P-AUTO-01}{}
\begin{sujettype}[P-AUTO-01 — Automatismes niveau épreuve]
\textbf{Type :} entraînement d'automatismes.  
\textbf{Durée conseillée :} 20 minutes.  
\textbf{Calculatrice interdite.}

Répondre sans justification.

\begin{enumerate}
  \item Calculer :
  \[
  \frac{3}{4}+\frac{5}{6}.
  \]
  \item Simplifier :
  \[
  \frac{12x^2}{18x}.
  \]
  \item Développer :
  \[
  (2x-3)^2.
  \]
  \item Factoriser :
  \[
  x^2-16.
  \]
  \item Résoudre :
  \[
  3x-7=11.
  \]
  \item Résoudre :
  \[
  -2x+5\leq 9.
  \]
  \item Calculer :
  \[
  f(2)\quad\text{pour}\quad f(x)=x^2-3x+1.
  \]
  \item Donner la dérivée de :
  \[
  f(x)=3x^2-5x+2.
  \]
  \item Calculer :
  \[
  \overrightarrow{AB}
  \quad\text{avec}\quad
  A(1;3),\ B(5;-2).
  \]
  \item Calculer :
  \[
  P(\overline A)
  \quad\text{si}\quad
  P(A)=0{,}37.
  \]
\end{enumerate}
\end{sujettype}

\begin{corrige}[P-AUTO-01]
\begin{multicols}{2}
\begin{enumerate}
  \item $\frac{19}{12}$.
  \item $\frac{2x}{3}$ si $x\neq0$.
  \item $4x^2-12x+9$.
  \item $(x-4)(x+4)$.
  \item $x=6$.
  \item $x\geq -2$.
  \item $f(2)=4-6+1=-1$.
  \item $f'(x)=6x-5$.
  \item $\overrightarrow{AB}(4;-5)$.
  \item $0{,}63$.
\end{enumerate}
\end{multicols}
\end{corrige}

% ------------------------------------------------------------
\hypertarget{P-SUJET-01}{}
\begin{sujettype}[P-SUJET-01 — Diagnostic première spécialité]
\textbf{Type :} sujet complet inventé.  
\textbf{Durée :} 2 h.  
\textbf{Calculatrice interdite.}

\subsection*{Partie A — Automatismes}

\begin{enumerate}
  \item Développer et réduire :
  \[
  A=(x-4)^2-2(x+1).
  \]
  \item Factoriser :
  \[
  B=9x^2-25.
  \]
  \item Résoudre :
  \[
  x^2-5x+6=0.
  \]
  \item Résoudre :
  \[
  x^2-4x+3\leq0.
  \]
\end{enumerate}

\subsection*{Partie B — Fonction}

On considère la fonction
\[
f(x)=x^3-3x^2-9x+5.
\]

\begin{enumerate}
  \item Calculer $f'(x)$.
  \item Factoriser $f'(x)$.
  \item Étudier le signe de $f'(x)$.
  \item Dresser le tableau de variations de $f$.
  \item Déterminer les extrema locaux de $f$.
\end{enumerate}

\subsection*{Partie C — Probabilités}

Dans une classe de 32 élèves, 18 suivent la spécialité mathématiques, 14 suivent la spécialité physique-chimie et 9 suivent les deux spécialités.

On choisit un élève au hasard.

\begin{enumerate}
  \item Définir deux événements adaptés.
  \item Calculer la probabilité que l'élève suive les deux spécialités.
  \item Calculer la probabilité que l'élève suive au moins une des deux spécialités.
  \item Calculer la probabilité que l'élève ne suive aucune des deux.
\end{enumerate}
\end{sujettype}

% ------------------------------------------------------------
\hypertarget{P-SUJET-02}{}
\begin{sujettype}[P-SUJET-02 — Second degré et paraboles]
\textbf{Type :} exercice long inventé.  
\textbf{Durée :} 1 h.  
\textbf{Chapitres :} second degré, signe, optimisation.

Une entreprise fabrique des objets. Le bénéfice, en centaines d'euros, est modélisé par :
\[
B(x)=-2x^2+24x-40,
\]
où $x$ désigne le nombre de centaines d'objets produits.

\begin{enumerate}
  \item Calculer $B(0)$, $B(2)$ et $B(10)$.
  \item Résoudre $B(x)=0$.
  \item Factoriser $B(x)$.
  \item Étudier le signe de $B(x)$.
  \item Interpréter l'intervalle sur lequel l'entreprise réalise un bénéfice positif.
  \item Mettre $B(x)$ sous forme canonique.
  \item Déterminer la production qui maximise le bénéfice.
  \item Donner le bénéfice maximal.
\end{enumerate}
\end{sujettype}

\begin{corrige}[P-SUJET-02]
On a
\[
B(x)=-2x^2+24x-40=-2(x^2-12x+20).
\]
L'équation $B(x)=0$ revient à :
\[
x^2-12x+20=0.
\]
Le discriminant vaut :
\[
\Delta=144-80=64.
\]
Les racines sont :
\[
x_1=2,\qquad x_2=10.
\]
Donc :
\[
B(x)=-2(x-2)(x-10).
\]
Comme le coefficient dominant est négatif, le bénéfice est positif entre les racines :
\[
B(x)>0\quad\text{pour}\quad x\in]2;10[.
\]
La forme canonique est :
\[
B(x)=-2(x-6)^2+32.
\]
Le bénéfice maximal vaut donc $32$ centaines d'euros, atteint pour $x=6$.
\end{corrige}

% ------------------------------------------------------------
\hypertarget{P-SUJET-03}{}
\begin{sujettype}[P-SUJET-03 — Fonction et lecture graphique]
\textbf{Type :} exercice long inventé.  
\textbf{Durée :} 1 h.  
\textbf{Chapitres :} fonctions, second degré, variations.

On considère la fonction
\[
f(x)=-x^2+6x-5.
\]

\begin{enumerate}
  \item Calculer $f(0)$, $f(1)$, $f(3)$ et $f(6)$.
  \item Résoudre $f(x)=0$.
  \item Mettre $f(x)$ sous forme canonique.
  \item En déduire le maximum de $f$.
  \item Dresser le tableau de variations de $f$ sur $\R$.
  \item Résoudre $f(x)\geq0$.
  \item Interpréter graphiquement l'ensemble des solutions.
\end{enumerate}
\end{sujettype}

% ------------------------------------------------------------
\hypertarget{P-SUJET-04}{}
\begin{sujettype}[P-SUJET-04 — Dérivation et tangentes]
\textbf{Type :} exercice long inventé.  
\textbf{Durée :} 1 h 15.  
\textbf{Chapitres :} dérivation, tangente, variations.

On considère la fonction
\[
f(x)=x^3-6x^2+9x+2.
\]

\begin{enumerate}
  \item Calculer $f'(x)$.
  \item Résoudre $f'(x)=0$.
  \item Étudier le signe de $f'(x)$.
  \item Dresser le tableau de variations de $f$.
  \item Calculer $f(1)$ et $f'(1)$.
  \item Écrire l'équation de la tangente à la courbe de $f$ au point d'abscisse $1$.
  \item Déterminer les extrema locaux de $f$.
\end{enumerate}
\end{sujettype}

\begin{corrige}[P-SUJET-04]
\[
f'(x)=3x^2-12x+9=3(x^2-4x+3)=3(x-1)(x-3).
\]
Ainsi $f'(x)$ est positif sur $]-\infty;1]$, négatif sur $[1;3]$, puis positif sur $[3;+\infty[$.  
Donc $f$ est croissante, puis décroissante, puis croissante.  
On a :
\[
f(1)=1-6+9+2=6,
\qquad
f'(1)=0.
\]
La tangente en $1$ est donc :
\[
y=6.
\]
\end{corrige}

% ------------------------------------------------------------
\hypertarget{P-SUJET-05}{}
\begin{sujettype}[P-SUJET-05 — Sujet complet type épreuve 1]
\textbf{Type :} sujet complet inventé.  
\textbf{Durée :} 2 h.  
\textbf{Calculatrice interdite.}

\subsection*{Exercice 1 — Automatismes}

\begin{enumerate}
  \item Calculer :
  \[
  \frac{2}{3}-\frac{5}{12}.
  \]
  \item Développer :
  \[
  (3x+2)^2.
  \]
  \item Factoriser :
  \[
  4x^2-12x.
  \]
  \item Résoudre :
  \[
  x^2-7x+12=0.
  \]
  \item Donner la dérivée de :
  \[
  f(x)=5x^3-2x^2+7.
  \]
\end{enumerate}

\subsection*{Exercice 2 — Étude de fonction}

Soit
\[
f(x)=2x^3-9x^2+12x-1.
\]

\begin{enumerate}
  \item Calculer $f'(x)$.
  \item Factoriser $f'(x)$.
  \item Étudier le signe de $f'(x)$.
  \item Dresser le tableau de variations.
  \item Déterminer les extrema locaux.
\end{enumerate}

\subsection*{Exercice 3 — Suites}

On définit la suite $(u_n)$ par :
\[
u_0=120,
\qquad
u_{n+1}=0{,}85u_n+18.
\]

\begin{enumerate}
  \item Calculer $u_1$ et $u_2$.
  \item Résoudre l'équation $L=0{,}85L+18$.
  \item On pose $v_n=u_n-L$. Montrer que $(v_n)$ est géométrique.
  \item En déduire une expression de $u_n$.
  \item Interpréter la limite de la suite.
\end{enumerate}
\end{sujettype}

% ------------------------------------------------------------
\hypertarget{P-SUJET-06}{}
\begin{sujettype}[P-SUJET-06 — Suites et modélisation]
\textbf{Type :} exercice long inventé.  
\textbf{Durée :} 1 h 15.  
\textbf{Chapitres :} suites, géométrique, seuil, algorithme.

Une chaîne YouTube compte initialement 1500 abonnés. Chaque mois, elle gagne 8 \% d'abonnés supplémentaires.

On note $u_n$ le nombre d'abonnés après $n$ mois.

\begin{enumerate}
  \item Donner $u_0$.
  \item Justifier que $u_{n+1}=1{,}08u_n$.
  \item Reconnaître la nature de la suite.
  \item Donner l'expression de $u_n$ en fonction de $n$.
  \item Calculer $u_6$ à l'unité près.
  \item Déterminer au bout de combien de mois la chaîne dépassera 3000 abonnés.
  \item Écrire un algorithme Python permettant de trouver ce seuil.
\end{enumerate}
\end{sujettype}

% ------------------------------------------------------------
\hypertarget{P-SUJET-07}{}
\begin{sujettype}[P-SUJET-07 — Probabilités dans un lycée]
\textbf{Type :} exercice long inventé.  
\textbf{Durée :} 1 h.  
\textbf{Chapitres :} probabilités, union, intersection, contraire.

Dans un lycée, 48 \% des élèves suivent la spécialité mathématiques, 36 \% suivent la spécialité physique-chimie et 22 \% suivent les deux spécialités.

On choisit un élève au hasard.

On note :
\[
M : \text{``l'élève suit la spécialité mathématiques''},
\]
\[
P : \text{``l'élève suit la spécialité physique-chimie''}.
\]

\begin{enumerate}
  \item Donner $P(M)$, $P(P)$ et $P(M\cap P)$.
  \item Calculer $P(M\cup P)$.
  \item Calculer la probabilité que l'élève ne suive aucune des deux spécialités.
  \item Calculer la probabilité que l'élève suive mathématiques mais pas physique-chimie.
  \item Parmi les élèves suivant mathématiques, calculer la proportion suivant aussi physique-chimie.
\end{enumerate}
\end{sujettype}

% ------------------------------------------------------------
\hypertarget{P-SUJET-08}{}
\begin{sujettype}[P-SUJET-08 — Test médical et probabilités conditionnelles]
\textbf{Type :} exercice long inventé.  
\textbf{Durée :} 1 h 15.  
\textbf{Chapitres :} arbre pondéré, probabilités conditionnelles.

Une maladie touche 3 \% d'une population.  
Un test est positif chez 98 \% des personnes malades et chez 6 \% des personnes non malades.

On note :
\[
M : \text{``la personne est malade''},
\qquad
T : \text{``le test est positif''}.
\]

\begin{enumerate}
  \item Donner $P(M)$ et $P(\overline M)$.
  \item Donner $P_M(T)$ et $P_{\overline M}(T)$.
  \item Construire un arbre pondéré.
  \item Calculer $P(M\cap T)$.
  \item Calculer $P(\overline M\cap T)$.
  \item En déduire $P(T)$.
  \item Calculer $P_T(M)$.
  \item Interpréter le résultat obtenu.
\end{enumerate}
\end{sujettype}

\begin{corrige}[P-SUJET-08]
\[
P(M)=0{,}03,
\qquad
P(\overline M)=0{,}97.
\]
\[
P_M(T)=0{,}98,
\qquad
P_{\overline M}(T)=0{,}06.
\]
Donc :
\[
P(M\cap T)=0{,}03\times0{,}98=0{,}0294.
\]
Et :
\[
P(\overline M\cap T)=0{,}97\times0{,}06=0{,}0582.
\]
Ainsi :
\[
P(T)=0{,}0294+0{,}0582=0{,}0876.
\]
Donc :
\[
P_T(M)=\frac{P(M\cap T)}{P(T)}
=\frac{0{,}0294}{0{,}0876}
\approx0{,}336.
\]
Même si le test est positif, la probabilité d'être malade est environ $33{,}6\%$.
\end{corrige}

% ------------------------------------------------------------
\hypertarget{P-SUJET-09}{}
\begin{sujettype}[P-SUJET-09 — Géométrie repérée]
\textbf{Type :} exercice long inventé.  
\textbf{Durée :} 1 h.  
\textbf{Chapitres :} coordonnées, vecteurs, produit scalaire.

Dans un repère orthonormé, on donne :
\[
A(1;2),\qquad B(5;4),\qquad C(3;8).
\]

\begin{enumerate}
  \item Calculer les coordonnées de $\overrightarrow{AB}$ et $\overrightarrow{AC}$.
  \item Calculer $AB$ et $AC$.
  \item Calculer le produit scalaire $\overrightarrow{AB}\cdot\overrightarrow{AC}$.
  \item Le triangle $ABC$ est-il rectangle en $A$ ?
  \item Calculer les coordonnées du milieu de $[BC]$.
  \item Déterminer une équation de la droite $(AB)$.
\end{enumerate}
\end{sujettype}

% ------------------------------------------------------------
\hypertarget{P-SUJET-10}{}
\begin{sujettype}[P-SUJET-10 — Sujet complet type épreuve 2]
\textbf{Type :} sujet complet inventé.  
\textbf{Durée :} 2 h.  
\textbf{Calculatrice interdite.}

\subsection*{Exercice 1 — Automatismes}

\begin{enumerate}
  \item Calculer :
  \[
  \frac{7}{10}+\frac{3}{5}.
  \]
  \item Résoudre :
  \[
  4x+3=2x-5.
  \]
  \item Factoriser :
  \[
  x^2+6x+9.
  \]
  \item Donner la dérivée de :
  \[
  f(x)=x^4-3x^2+1.
  \]
  \item Calculer :
  \[
  P(\overline A)
  \quad\text{si}\quad
  P(A)=\frac{7}{12}.
  \]
\end{enumerate}

\subsection*{Exercice 2 — Second degré}

Soit
\[
f(x)=x^2-8x+12.
\]

\begin{enumerate}
  \item Résoudre $f(x)=0$.
  \item Factoriser $f(x)$.
  \item Étudier le signe de $f(x)$.
  \item Mettre $f(x)$ sous forme canonique.
  \item En déduire le minimum de $f$.
\end{enumerate}

\subsection*{Exercice 3 — Probabilités conditionnelles}

Une entreprise fabrique des pièces à l'aide de deux machines.  
La machine $A$ produit 70 \% des pièces et la machine $B$ en produit 30 \%.  
La machine $A$ produit 2 \% de pièces défectueuses.  
La machine $B$ produit 5 \% de pièces défectueuses.

On choisit une pièce au hasard.

\begin{enumerate}
  \item Construire un arbre pondéré.
  \item Calculer la probabilité que la pièce vienne de $A$ et soit défectueuse.
  \item Calculer la probabilité que la pièce soit défectueuse.
  \item Une pièce est défectueuse. Calculer la probabilité qu'elle provienne de $B$.
\end{enumerate}
\end{sujettype}

% ------------------------------------------------------------
\hypertarget{P-SUJET-11}{}
\begin{sujettype}[P-SUJET-11 — Trigonométrie]
\textbf{Type :} exercice long inventé.  
\textbf{Durée :} 1 h.  
\textbf{Chapitres :} cercle trigonométrique, sinus, cosinus.

\begin{enumerate}
  \item Donner les valeurs exactes de :
  \[
  \cos(0),\quad \sin(0),\quad \cos\left(\frac{\pi}{3}\right),\quad
  \sin\left(\frac{\pi}{3}\right).
  \]
  \item Donner les valeurs exactes de :
  \[
  \cos\left(\frac{\pi}{4}\right),\quad
  \sin\left(\frac{\pi}{4}\right),\quad
  \cos\left(\frac{\pi}{6}\right),\quad
  \sin\left(\frac{\pi}{6}\right).
  \]
  \item Résoudre dans $[0;2\pi]$ :
  \[
  \cos(x)=\frac12.
  \]
  \item Résoudre dans $[0;2\pi]$ :
  \[
  \sin(x)=\frac{\sqrt3}{2}.
  \]
  \item Placer les solutions sur le cercle trigonométrique.
\end{enumerate}
\end{sujettype}

% ------------------------------------------------------------
\hypertarget{P-SUJET-12}{}
\begin{sujettype}[P-SUJET-12 — Algorithmique et suites]
\textbf{Type :} exercice long inventé.  
\textbf{Durée :} 1 h.  
\textbf{Chapitres :} Python, suites, seuil.

On considère la suite définie par :
\[
u_0=50,
\qquad
u_{n+1}=1{,}12u_n+8.
\]

\begin{enumerate}
  \item Calculer $u_1$ et $u_2$.
  \item Recopier et compléter le programme suivant pour calculer $u_{20}$ :
\begin{verbatim}
u = 50
for n in range(...):
    u = ...
print(u)
\end{verbatim}
  \item Écrire un programme qui détermine le premier rang $n$ tel que $u_n>200$.
  \item Expliquer le rôle de la boucle \texttt{while}.
  \item Interpréter le résultat dans un contexte de croissance.
\end{enumerate}
\end{sujettype}

% ------------------------------------------------------------
\hypertarget{P-SUJET-13}{}
\begin{sujettype}[P-SUJET-13 — Sujet complet type épreuve 3]
\textbf{Type :} sujet complet inventé.  
\textbf{Durée :} 2 h.  
\textbf{Calculatrice interdite.}

\subsection*{Exercice 1 — Automatismes}

\begin{enumerate}
  \item Calculer :
  \[
  \frac{5}{6}\times \frac{9}{10}.
  \]
  \item Développer :
  \[
  (x-7)(x+7).
  \]
  \item Résoudre :
  \[
  -3x+4\geq 10.
  \]
  \item Donner la dérivée de :
  \[
  f(x)=2x^3-5x+4.
  \]
  \item Calculer le milieu de $A(2;-1)$ et $B(6;5)$.
\end{enumerate}

\subsection*{Exercice 2 — Fonction et dérivée}

On considère :
\[
f(x)=-x^3+3x^2+9x-2.
\]

\begin{enumerate}
  \item Calculer $f'(x)$.
  \item Résoudre $f'(x)=0$.
  \item Étudier le signe de $f'(x)$.
  \item Dresser le tableau de variations de $f$.
  \item Déterminer les extrema locaux.
\end{enumerate}

\subsection*{Exercice 3 — Suite}

Une population de bactéries augmente de 15 \% par heure.  
Au départ, elle contient 800 bactéries.

\begin{enumerate}
  \item Modéliser la situation par une suite.
  \item Donner l'expression explicite de cette suite.
  \item Déterminer le nombre de bactéries après 5 heures.
  \item Écrire un algorithme qui détermine au bout de combien d'heures la population dépasse 5000.
\end{enumerate}

\subsection*{Exercice 4 — Géométrie}

Dans un repère orthonormé, on donne :
\[
A(0;1),\quad B(4;3),\quad C(2;7).
\]

\begin{enumerate}
  \item Calculer $\overrightarrow{AB}$ et $\overrightarrow{AC}$.
  \item Calculer $\overrightarrow{AB}\cdot\overrightarrow{AC}$.
  \item Le triangle $ABC$ est-il rectangle en $A$ ?
  \item Déterminer une équation de la droite $(AB)$.
\end{enumerate}
\end{sujettype}

% ============================================================
\section{Sujets plus difficiles que l'épreuve}
% ============================================================

\begin{alerte}[utilisation]
Ces sujets sont plus difficiles que l'épreuve. Ils sont utiles pour un très bon élève qui vise une excellente note ou une poursuite en prépa, mais ils ne doivent pas remplacer les automatismes et les sujets zéro.
\end{alerte}

\hypertarget{P-DIFF-01}{}
\begin{difficile}[P-DIFF-01 — Fonction avec paramètre]
\textbf{Durée :} 1 h 30.  
\textbf{Niveau :} difficile première / transition terminale.

Pour tout réel $a$, on considère :
\[
f_a(x)=x^3-3ax^2+3x+1.
\]

\begin{enumerate}
  \item Calculer $f_a'(x)$.
  \item Étudier selon $a$ le nombre de solutions de $f_a'(x)=0$.
  \item Pour $a=2$, dresser le tableau de variations de $f_a$.
  \item Pour $a=1$, étudier les variations de $f_a$.
  \item Interpréter le rôle du paramètre $a$.
\end{enumerate}
\end{difficile}

\hypertarget{P-DIFF-02}{}
\begin{difficile}[P-DIFF-02 — Suite affine et changement de variable]
\textbf{Durée :} 1 h 30.  
\textbf{Niveau :} difficile première.

On définit :
\[
u_0=10,
\qquad
u_{n+1}=0{,}7u_n+12.
\]

\begin{enumerate}
  \item Calculer les quatre premiers termes.
  \item Résoudre $L=0{,}7L+12$.
  \item Poser $v_n=u_n-L$.
  \item Montrer que $(v_n)$ est géométrique.
  \item Donner l'expression de $u_n$.
  \item Déterminer la limite de $(u_n)$.
  \item Écrire un programme de seuil.
\end{enumerate}
\end{difficile}

\hypertarget{P-DIFF-03}{}
\begin{difficile}[P-DIFF-03 — Probabilités et décision]
\textbf{Durée :} 1 h 30.  
\textbf{Niveau :} difficile première.

Une entreprise affirme que 90 \% de ses colis arrivent à l'heure.  
On observe 50 colis.

\begin{enumerate}
  \item Proposer une variable aléatoire adaptée.
  \item Donner sa loi.
  \item Calculer son espérance.
  \item Interpréter l'événement $X\leq40$.
  \item Expliquer pourquoi une observation de 40 colis à l'heure peut remettre en question l'affirmation.
  \item Proposer une méthode de décision statistique simple.
\end{enumerate}
\end{difficile}

% ============================================================
\section{Carnet d'erreurs}
% ============================================================

\begin{methode}[comment corriger activement]
Après chaque sujet, remplir le tableau ci-dessous.  
Le but n'est pas seulement de connaître sa note, mais de comprendre précisément les erreurs.
\end{methode}

\footnotesize
\begin{longtable}{|p{1.8cm}|p{2.5cm}|p{2cm}|p{4cm}|p{4cm}|p{2cm}|}
\hline
\textbf{Date} & \textbf{Sujet} & \textbf{Question} & \textbf{Erreur commise} & \textbf{Correction à retenir} & \textbf{Refait ?}\\
\hline
 & & & & & \\
\hline
 & & & & & \\
\hline
 & & & & & \\
\hline
 & & & & & \\
\hline
 & & & & & \\
\hline
 & & & & & \\
\hline
 & & & & & \\
\hline
 & & & & & \\
\hline
\end{longtable}
\normalsize

% ============================================================
\section{Check-list finale avant l'épreuve}
% ============================================================

\begin{alerte}[à ne pas faire]
Dans les dernières 48 heures, ne plus commencer un gros nouveau chapitre.  
Il faut consolider, relire, refaire les erreurs, dormir correctement et arriver lucide.
\end{alerte}

\begin{multicols}{2}
\begin{itemize}
  \item[\casevide] Je sais calculer sans calculatrice.
  \item[\casevide] Je sais développer et factoriser.
  \item[\casevide] Je sais résoudre une équation du second degré.
  \item[\casevide] Je sais étudier le signe d'un trinôme.
  \item[\casevide] Je sais calculer une dérivée.
  \item[\casevide] Je sais faire un tableau de variations.
  \item[\casevide] Je sais écrire une tangente.
  \item[\casevide] Je sais manipuler une suite arithmétique.
  \item[\casevide] Je sais manipuler une suite géométrique.
  \item[\casevide] Je sais utiliser un arbre pondéré.
  \item[\casevide] Je sais calculer une probabilité conditionnelle.
  \item[\casevide] Je sais calculer des coordonnées de vecteurs.
  \item[\casevide] Je sais utiliser le produit scalaire.
  \item[\casevide] Je sais lire un algorithme Python.
  \item[\casevide] J'ai fait au moins deux sujets complets en 2 h.
  \item[\casevide] J'ai relu mon carnet d'erreurs.
\end{itemize}
\end{multicols}

% ============================================================
\section{Conseils pour le jour de l'épreuve}
% ============================================================

\begin{methode}[pendant l'épreuve]
\begin{enumerate}
  \item Commencer par respirer et lire le sujet en entier.
  \item Faire les automatismes rapidement mais proprement.
  \item Ne pas rester bloqué sur une question.
  \item Dans les exercices, rédiger les conclusions.
  \item Encadrer les résultats importants.
  \item Vérifier les signes et les domaines.
  \item Garder au moins 10 minutes de relecture.
\end{enumerate}
\end{methode}

\begin{alerte}[erreurs classiques qui coûtent cher]
\begin{multicols}{2}
\begin{itemize}
  \item Oublier de changer le sens d'une inégalité.
  \item Confondre racines et signe d'un trinôme.
  \item Mal factoriser une dérivée.
  \item Oublier la phrase de conclusion.
  \item Confondre $P(A\cap B)$ et $P_A(B)$.
  \item Mal lire les coordonnées d'un vecteur.
  \item Écrire une tangente avec $f'(x)$ au lieu de $f'(a)$.
  \item Utiliser la calculatrice alors qu'elle est interdite.
\end{itemize}
\end{multicols}
\end{alerte}

\vfill

\begin{center}
\Large
\textbf{Objectif : une copie propre, rapide, précise et complète.}\\[1mm]
\textbf{La meilleure note vient d'abord des points faciles parfaitement sécurisés.}
\end{center}

\end{document}